
\documentclass[11pt]{article}
\usepackage[T1]{fontenc}
\usepackage[utf8]{inputenc}
\usepackage{lmodern,microtype,amsmath,amssymb,amsthm,mathtools}
\usepackage[a4paper,margin=27mm]{geometry}
\usepackage{graphicx,xcolor,booktabs,longtable,array,enumitem,tikz}
\usetikzlibrary{arrows.meta,calc}
\usepackage[unicode,colorlinks=true,linkcolor=blue!45!black,citecolor=green!35!black,urlcolor=blue!55!black]{hyperref}
\usepackage{xurl}
\setlength{\emergencystretch}{2em}
\setlist{itemsep=3pt,topsep=5pt}
\numberwithin{equation}{section}
\newtheorem{theorem}{Theorem}
\newtheorem{lemma}[theorem]{Lemma}
\newtheorem{proposition}[theorem]{Proposition}
\newtheorem{foundation}{Companion theorem}
\providecommand{\tightlist}{\setlength{\itemsep}{0pt}\setlength{\parskip}{0pt}}
\title{Volume limits for long alternating-generator three-braids}
\author{}
\date{1 November 2026}
\begin{document}
\maketitle
\begin{abstract}
Fix an even number \(m\ge6\) of syllables in a closed three-strand braid,
with generators alternating between \(\sigma_2\) and \(\sigma_1\).
For independently sufficiently long signed powers, whenever the closure
is a knot, its reduced Kashaev invariant has exponential rate equal to
the hyperbolic volume of the knot complement. The limit includes both
parities and composite levels. We identify the original crossing disks
with an explicitly matched bipyramid and compute its marked packing.
Three shared choices of face weights give a uniform \(7/100\) loss for
every complementary source sector, for every fixed even \(m\).
The full all-size estimate and the fixed rational constant bounds are
proved here. Exact finite source, local cofactor, and nonzero geometric
filling results are stated explicitly as companion inputs; their
analytic and geometric proofs are not reproduced in this paper.
\end{abstract}

The volume question originates with Kashaev~\cite{Kashaev} and its
colored-Jones formulation with Murakami and Murakami~\cite{Murakami}.
Here every diagram and every signed power is fixed before the quantum
level tends to infinity. The phrase \emph{alternating generators}
describes the braid word, not the crossing signs of its projection.
Geometric volume estimates for closed three-braids appear in
Futer--Kalfagianni--Purcell~\cite{FKP}; no geometric bound from that work
is used as a substitute for the root-of-unity asymptotic below.

The proof separates its three companion inputs from its new comparison.
Companion Theorem~\ref{found:finite} identifies the original invariant
with a shared coefficient and a bounded finite filling functional.
Companion Theorem~\ref{found:envelopes} supplies phase-uniform local
rates. Companion Theorem~\ref{found:period} evaluates the nonzero central
contribution and identifies its action with the actual filled volume.
The main work here is a shared, global bound for all the other terms;
this is what permits the central contribution to determine the full
invariant. A bound for isolated factors with unrelated face weights
would not suffice.

\section{The sphere and the
theorem}\label{braid:sec-1}

Throughout, fix an even integer \(m\ge6\). In the diagram source the
quantum level is \(N\ge3\), and \(d=N-1\). Use poles \(N_0,S_0\),
rim vertices \(v_0,\ldots,v_{m-1}\), and the
oriented triangles \begin{equation}\label{braid:eq-1}
(N_0,v_j,v_{j+1}),\qquad(S_0,v_{j+1},v_j),
 \qquad j\in\mathbb Z/m.
\end{equation} Match north triangles across \(N_0v_j\) for odd \(j\), and south
triangles across \(S_0v_j\) for even \(j\). Every triangle belongs to
exactly one matched pair. This is a simple triangulated sphere with
\(m\) matched pairs and \(m+2\) white vertices.

The cyclic quadrilateral order is \((x,b,y,c)\); the homogeneous source
argument order is \((x,b,c,y)\). The original frames are \begin{equation}\label{braid:eq-2}
\begin{array}{c|c|c}
 & (x,b,c,y)&\text{matching edge}\\ \hline
 j\text{ odd}&(N_0,v_{j-1},v_{j+1},v_j)&N_0v_j\\
 j\text{ even}&(S_0,v_{j+1},v_{j-1},v_j)&S_0v_j .
\end{array}
\end{equation} The reversed south middle order is essential. Mark the unmatched edge
\(v_0v_1\).

Write \(R\) for the positive half-twist of the two strands in the
original disk frame; \(R^n\) is its signed integral power. Add one
crossing circle around each displayed box. Removing its full twists,
while retaining the chosen signed half twist, gives the augmented
exterior \(Y_m^\epsilon\) with these original disk markings.
Retain a signed half twist \(\epsilon_j\in\{-1,0,1\}\) at each original
crossing disk. Write \(\mu_j\) for its meridian and \(\lambda_j\) for
the disk-framed primitive longitude, with the orientation convention of
the original two-strand crossing. Filling kills \(\mu_j+p_j\lambda_j\).
In this basis, filling parameter
\(p_j\in\mathbb Z\) inserts the power \begin{equation}\label{braid:eq-3}
n_j=\epsilon_j-2p_j.
\end{equation} Let the braid generators use the same positive crossing as \(R\), and
write factors in the order encountered by the strands.

\begin{theorem}\label{thm:braids} The filled diagram is the ordinary closure of \begin{equation}\label{braid:eq-4}
\beta_m(\mathbf n)=
 \sigma_2^{n_0}\sigma_1^{n_1}\sigma_2^{n_2}\sigma_1^{n_3}
       \cdots\sigma_2^{n_{m-2}}\sigma_1^{n_{m-1}}.
\end{equation} Denote this closure by \(K_{m,\mathbf n}\). The \(m\)
crossing disks are the disks surrounding those syllables.
The braid axis is not a link component. The main link is a knot exactly
when the permutation of \eqref{braid:eq-4} is a 3-cycle. In particular, exactly one odd
north exponent and one odd south exponent, with every other exponent
even, gives a knot for every \(m\).

For each fixed \(m\), the estimates below and the companion filling
results give a finite \(B_m\) such that, whenever the fixed integers
satisfy \(\min_j|n_j|>B_m\) and the knot test, the closure is hyperbolic
and \begin{equation}\label{braid:eq-5}
\lim_{N\to\infty}\frac{2\pi}{N}
       \log|K_N(\widehat{\beta_m(\mathbf n)})|
  =\operatorname{Vol}(S^3\setminus\widehat{\beta_m(\mathbf n)}).
\end{equation}
\end{theorem}

Here \(K_N\) is the zero-framed reduced Kashaev normalization with
unknot value one. The signs and relative sizes of the fixed powers are
unrestricted. There is no uniformity in growing \(m\), no growing-twist
assertion, and no moving-color conclusion.

\section{Exact strand maps and the three-strand
braid}\label{braid:sec-2}

Number the sides of a cyclic quadrilateral by \begin{equation}\label{braid:eq-6}
0=xb,\quad1=by,\quad2=yc,\quad3=cx.
\end{equation} These are the native cut corners \((C_1,C_2,C_3,C_0)\), respectively.
The literal pure shaded map \(C_1\mapsto C_0,\ C_2\mapsto C_3\) gives
main-strand pairs \((0,3),(1,2)\). Either signed odd half twist gives
\((0,2),(1,3)\). The shaded grouping of corners is a different pairing
and is not used to count main components.

Embed the sphere with \(N_0\) inside the equator and \(S_0\) outside.
The outside strips are exactly these:

\begin{itemize}
\tightlist
\item
  the edge \(v_jv_{j+1}\) joins boxes \(j,j+1\);
\item
  an unmatched north spoke, at even rim index \(j\), joins odd boxes
  \(j-1,j+1\);
\item
  an unmatched south spoke, at odd rim index \(j\), joins even boxes
  \(j-1,j+1\).
\end{itemize}

Under the pure pairing they form three nested planar circles: north,
equator, south. There are no other strips. Straighten these three
circles into concentric tracks while retaining the disjoint disk
neighborhoods. Increasing angular coordinate is braid time; label the
tracks from inside to outside by \(1,2,3\).

The forward endpoints at each box are \begin{equation}\label{braid:eq-7}
\begin{array}{c|c|c}
 &\text{input, inner then outer}
 &\text{output, inner then outer}\\ \hline
 j\text{ odd}
 &(N_0v_{j-1},v_{j-1}v_j)
 &(N_0v_{j+1},v_jv_{j+1})\\
 j\text{ even}
 &(v_{j-1}v_j,S_0v_{j-1})
 &(v_jv_{j+1},S_0v_{j+1}).
\end{array}
\end{equation} North boxes therefore act on tracks \(1,2\); south boxes act on
tracks \(2,3\). The unused track simply passes the box through its
outside strip.

The north source frame already has increasing angular time and the
inside track on its left. At a south box, both the time direction and
the transverse order are reversed. Thus the change from its source frame
to \eqref{braid:eq-7} is a planar half-turn. It is not a reflection or a one-head
orientation reversal.

For a single crossing a planar half-turn preserves each diagonal pair of
endpoints and therefore preserves which diagonal is over. It retains the
signed crossing \(R\). For \(R^n\) the order of the identical signed
crossings is reversed, and the exponent is still \(n\). Reversing both
strand orientations also preserves crossing writhe. This is a literal
tangle isotopy with the ribbon normal transported; it is not the
assertion \(T^2=I\) for typed quantum turns.

Cut the annulus before box 0. Formula \eqref{braid:eq-7} gives \eqref{braid:eq-4} with no sign
inversion at the south boxes. Closing the cut tracks gives the ordinary
braid closure. The half-turn and straightening carry the original
crossing circles and their disk collars along, so the identification
preserves the designated twist regions rather than just the unmarked
link.

Components of a braid closure are the cycles of its strand permutation.
With one odd box of each kind the permutation is a product of the two
distinct transpositions \((12),(23)\), hence a 3-cycle, in either order.
All even exponents give three components. All odd exponents give \begin{equation}\label{braid:eq-8}
\#\text{main components}=\gcd(3,m/2).
\end{equation} These statements hold for every \(m\) by \eqref{braid:eq-7}, not by a finite
experiment. More general parity patterns are tested exactly by the
permutation of \eqref{braid:eq-4}.

The phrase ``alternating generators'' does not assert an alternating
link diagram. For example, assigning opposite signs to the two generator
types gives the usual alternating braid pattern; arbitrary signs need
not.


\begin{figure}[htbp]
\centering
\resizebox{\textwidth}{!}{\definecolor{ink}{HTML}{183444}
\definecolor{north}{HTML}{1B6780}
\definecolor{south}{HTML}{B56722}
\begin{tikzpicture}[x=1cm,y=1cm,font=\small,text=ink,line cap=round,line join=round,
 edge/.style={draw=ink!45,line width=.7pt},
 markN/.style={draw=north,line width=2.1pt},
 markS/.style={draw=south,line width=2.1pt},
 dot/.style={circle,fill=ink,inner sep=1.3pt},
 move/.style={-{Latex[length=2mm]},draw=ink,line width=.9pt}]
\node[font=\bfseries] at (-4.15,3.05) {The matched white sphere ($m=6$)};
% Opposite orientations on the two equatorial boundary circles.
\foreach \name/\cx/\dir/\pole in {N/-6.35/1/N_0,S/-2.05/-1/S_0}{
 \begin{scope}[shift={(\cx,0)}]
  \fill[ink!3] (0,0) circle (1.65);
  \draw[edge] (0,0) circle (1.65);
  \foreach \j in {0,...,5}{
   \pgfmathsetmacro{\a}{90+\dir*60*\j}
   \coordinate (\name\j) at (\a:1.65);
   \draw[edge] (0,0)--(\name\j);
   \node[dot] at (\name\j) {};
   \node at (\a:1.97) {$v_{\j}$};
  }
  \ifnum\dir=1
   \foreach \j in {1,3,5}{
    \draw[markN] (0,0)--(\name\j);
    \path (0,0)--(\name\j) node[pos=.59,fill=white,inner sep=1.5pt,
       text=north,font=\scriptsize] {$j=\j$};
   }
  \else
   \foreach \j in {0,2,4}{
    \draw[markS] (0,0)--(\name\j);
    \path (0,0)--(\name\j) node[pos=.59,fill=white,inner sep=1.5pt,
       text=south,font=\scriptsize] {$j=\j$};
   }
  \fi
  \node[dot] at (0,0) {};
  \node[fill=white,inner sep=1.5pt] at (.29,-.18) {$\pole$};
 \end{scope}
}
\node[align=center,text=north] at (-6.35,-2.45) {North: match $N_0v_j$\\for odd $j$};
\node[align=center,text=south] at (-2.05,-2.45) {South: match $S_0v_j$\\for even $j$};
\node[align=center,font=\footnotesize] at (-4.2,-3.24)
 {Glue equal rim labels; the two boundary orders are opposite.};
\draw[move] (.15,0)--(1.15,0);
% The tracks are closed; a box replaces exactly its two adjacent track arcs.
\begin{scope}[shift={(4.55,0)}]
 \node[font=\bfseries] at (0,3.05) {The same disks on three annular tracks};
 \foreach \r in {1,1.75,2.5}{\draw[draw=ink,line width=1.1pt] (0,0) circle (\r);}
 \foreach \j in {0,...,5}{
  \pgfmathsetmacro{\a}{60*\j}
  \pgfmathsetmacro{\textangle}{\a>90 && \a<270 ? \a-180 : \a}
  \pgfmathtruncatemacro{\odd}{mod(\j,2)}
  \ifnum\odd=1
   \def\rc{1.375}\def\col{north}
  \else
   \def\rc{2.125}\def\col{south}
  \fi
  \begin{scope}[rotate=\a]
   \draw[draw=\col,fill=white,line width=1pt,rounded corners=1.5pt]
     (\rc-.54,-.25) rectangle (\rc+.54,.25);
   \node[font=\footnotesize,text=\col,rotate=\textangle] at (\rc,0) {$R^{n_{\j}}$};
  \end{scope}
 }
 \draw[move] (14:2.82) arc[start angle=14,end angle=42,radius=2.82];
 \node[font=\scriptsize,align=left] at (3.12,1.64) {braid\\time};
 \foreach \r/\i in {1/1,1.75/2,2.5/3}{
  \node[fill=white,inner sep=1pt,font=\footnotesize] at (225:\r) {$\i$};
 }
 \node[align=center,font=\footnotesize] at (0,-3.24)
  {Odd $j$: tracks $1,2$.\qquad Even $j$: tracks $2,3$.};
\end{scope}
\node[align=center,text width=16.3cm] at (-.75,-4.18)
 {$\displaystyle
 \beta_6=\sigma_2^{n_0}\sigma_1^{n_1}\sigma_2^{n_2}
         \sigma_1^{n_3}\sigma_2^{n_4}\sigma_1^{n_5}.$};
\end{tikzpicture}
}
\caption{The case \(m=6\). Equal rim labels on the two hemispheres are
identified with opposite boundary orientations. A thick matching spoke
pairs the two neighboring triangles and specifies one crossing disk.
The right-hand picture uses those same disks: odd indices act on the
inner two tracks, even indices on the outer two. Each box contains its
original signed power. The picture illustrates the endpoint table
\eqref{braid:eq-7}, whose proof applies to every even \(m\ge6\).}
\label{fig:matching}
\end{figure}

\section{Complete packing and the marked
curvatures}\label{braid:sec-3}

Put \begin{equation}\label{braid:eq-9}
\theta=\pi/m,\qquad a=\csc\theta.
\end{equation} In the complex projection plane take the north disk \(|z|\le a-1\),
the south \textbf{exterior} disk \(|z|\ge a+1\), and the unit disks \begin{equation}\label{braid:eq-10}
v_j:\quad |z-ae^{2ij\theta}|\le1.
\end{equation} Every rim disk is tangent to both pole disks. Rim centers at cyclic
separation \(r\) have distance \(2a\sin(r\theta)\), which is 2 exactly
for adjacent indices and strictly greater than 2 for every nonadjacent
pair. The two pole disks are disjoint. Thus the contact graph, with its
exterior south region, is exactly \eqref{braid:eq-1}; no extra edge is introduced.

The marked contact of \(v_0,v_1\) is \[
                  p=a\cos\theta\,e^{i\theta}.
\] Use the holomorphic Möbius transformation \begin{equation}\label{braid:eq-11}
w=\frac{2}{z-p}.
\end{equation} A further rotation or translation puts the two protected circles into
horizontal strip position without changing radii.

Under \(z\mapsto1/(z-p)\), a circle of center \(c\) and radius \(r\) has
new radius \(r/\bigl||c-p|^2-r^2\bigr|\). A zero denominator gives a
line. For the north and south boundaries the denominators before taking
absolute value are \(2(a-1)\) and \(-2(a+1)\). Both circles have
inversion radius \(1/2\), and their oriented disks become bounded
because \(p\) is outside both original disks. The factor 2 in \eqref{braid:eq-11} gives
both new radii one.

For rim circle \(j\), \[
\begin{split}
 |ae^{2ij\theta}-p|^2-1
 &=2a^2\cos\theta
       \{\cos\theta-\cos((2j-1)\theta)\}\\
 &=4a^2\cos\theta\sin(j\theta)\sin((j-1)\theta).
\end{split}
\] It vanishes exactly at \(j=0,1\) and is positive elsewhere. Therefore
the complete marked curvatures, meaning reciprocals of Euclidean circle
radii, are \begin{equation}\label{braid:eq-12}
L_{N_0}=L_{S_0}=1,\qquad
 k_j=L_{v_j}
 =\frac{\cos\theta}{\sin^2\theta}
       \{\cos\theta-\cos((2j-1)\theta)\}.
\end{equation} In particular, \begin{equation}\label{braid:eq-13}
k_0=k_1=0,\quad
 k_j>0\ (2\le j\le m-1),\quad
 k_2=k_{m-1}=4\cos^2\theta.
\end{equation} The two protected weights are exact zeros. They are never replaced by
small positive curvatures.

The strip width is 2, since the radius-one north circle is tangent to
both parallel lines. If another source convention uses width 1, a common
rescaling multiplies all ordinary curvatures by 2. The action is
unchanged by common scaling because total numerator degree equals total
ordinary target. Formula \eqref{braid:eq-12} is the normalization with pole curvatures
one.

The complete-geometry assertion used below has two ingredients. First,
the simple-sphere construction of \cite[Sections 2--3]{Packing} turns
this packing into the labeled right-angled ideal polyhedron: white
faces correspond to the displayed circles and shaded triangles to the
interstices. Its hypotheses hold because the contact graph just
computed is the simple triangulated sphere \eqref{braid:eq-1}.
Second, the original cut is determined by the matching. Cutting the
annular diagram along the projection sphere and the original crossing
half disks gives two balls with these labeled faces. White faces pair
across the projection copies; shaded faces pair according to
\eqref{braid:eq-2}, with the retained half twist exchanging their copies
and punctures when prescribed. These are precisely the inverse-cut
identifications of \cite[Section 2]{MorseGeometry}. Thus that construction
applies to \(Y_m^\epsilon\), with its original disks and main strands.
No change to a different shaded-face matching is made.

\section{The exact source and what must be
bounded}\label{source:sec-2}


\begin{foundation}[Exact source and finite filling norm]\label{found:finite}
For the original cut of a fixed connected marked simple matched sphere
at \(N\ge3\), let \(d=N-1\), let \(\Omega\) be the complete-color two-strand Casimir,
and form its polynomial numerator
\(V_N(\sigma;\Omega)=(2T_N(\sigma/2)+2)(\sigma-\Omega)^{-1}\).
The homogeneous face polynomial \(H_V(x,b,c,y;\sigma)\) is the Fourier
transform of this original matrix numerator in the four boundary
states, with the turn and crossing frames retained. It has total face
degree \(d\) and spectral degree at most \(d\).
Here \(T_N(\cos\phi)=\cos(N\phi)\). The full diagram source is its shared coefficient with target \(d\)
at every ordinary face and target zero at the two faces incident to
the marked strand edge. The raw filled evaluation is the tensor product
of the original signed-power functionals \(\mathcal L_N^{n_e}\) applied
to that polynomial, and differs from the zero-framed reduced invariant
only by its prescribed unit-modulus character. For every complex
polynomial \(H\) of degree at most \(d\),
\[
 |\mathcal L_N^n H|\le (1+|n|/2)\max_{-2\le\sigma\le2}|H(\sigma)|,
 \qquad n\in\mathbb Z.
\]
The scalar polynomial norm also holds at \(N=2\); the diagram source
is used only at \(N\ge3\). These statements hold for both level parities, with first normal data
at every double spectral root and the value at the odd simple root.
\end{foundation}

This is a companion result, not a theorem proved anew here.
The exact marked coefficient, signed-power evaluation and Kashaev
normalization are proved in \cite[Theorem 9.1, Corollary 9.2 and
Theorem 10.1]{FillingPaper}; its Theorem 1.1 proves the scalar norm.
That paper includes the local Fourier calculation, the shared face
grid and the actual cups, caps and single marked cut. We check the
diagram hypotheses on the explicit sphere below.
The local decomposition into a central
term and the two one-sided recurrence pieces is the one of
\cite[Sections 1--3]{Cofactors}; no positive replacement of the full
matrix numerator is made.

Put \(q=e^{\pi i/N}\), \(d=N-1\), and \begin{equation}\label{source:eq-A6}
\chi_N(z)=2T_N(z/2)+2,
 \qquad V_N(z;\Omega)=\chi_N(z)(z-\Omega)^{-1}.
\end{equation} Polynomial division defines the degree-at-most-\(d\) matrix
polynomial in \eqref{source:eq-A6}, including at every spectral root. It is the original
Casimir numerator, not the adjugate of an unrelated matrix. Let
\(H_V(x,b,c,y;z)\) be its original homogeneous four-face polynomial.

Writing the pole variables as \(A,B\) and the rim variables as \(z_j\),
the full shared spectral polynomial is exactly \begin{equation}\label{source:eq-A7}
\begin{split}
 S_N(\boldsymbol\sigma)
 ={}&[A^dB^d\prod_{j=2}^{m-1}z_j^d]\,
 \left.\prod_{j\,\mathrm{odd}}
 H_V(A,z_{j-1},z_{j+1},z_j;\sigma_j)
 \prod_{j\,\mathrm{even}}
 H_V(B,z_{j+1},z_{j-1},z_j;\sigma_j)
 \right|_{z_0=z_1=0}.
\end{split}
\end{equation} Every local factor has total face degree \(d\) and spectral degree at
most \(d\). To define its sectors, put \(h=\lfloor N/2\rfloor\).
An original decreasing matrix path from row \(j\) to row \(i\) has
middle monomial \(b^ic^{d-j}\). At even \(N=2h\), it is central when
\(i\le h-1\) and \(j\ge h\); at odd \(N=2h+1\), it is central when
\(i\le h\) and \(j\ge h\). The left sector has \(j<h\), and the
right sector consists of the remaining paths entirely above the central
cut. These three sets partition the original numerator paths. Let
\(S_{N,\mathrm C}\) select the central part at every box by polynomial
linearity, and set \begin{equation}\label{source:eq-A8}
D_N=S_N-S_{N,\mathrm C}.
\end{equation} The task is a strict real-cube bound on \eqref{source:eq-A8}. The complete drilled
rate of \(S_N\) alone would not give the smaller fixed-filled rate.

The general Morse compiler proves \eqref{source:eq-A7} on this diagram. In its proof
each actual box has factor \(N^{-1}H_V/(bc)^{d/2}\); the physical cups,
caps and unrestricted mark leave \(N^{-m-1}\). The primitive
face-potential lattice has \(m+1\) independent \(N\)-point grids,
canceling that factor. Endpoint support is imposed before specializing
those grids. The total numerator degree \(md\) forces the unique
nonnegative alias: the two marked targets are zero and every other
target is \(d\). This retains all repeated-node normal channels.

At even levels a literal integrated parity cochain is \begin{equation}\label{source:eq-A9}
t(v_j)=j\pmod2,\qquad t(N_0)=1,\qquad t(S_0)=0.
\end{equation} It agrees across matched spokes and differs across each strand edge;
in particular the two protected endpoints differ. The integrated coset
character at the unique supported target is one. No extra parity or
power of \(N\) multiplies \eqref{source:eq-A7}.

The raw filled contraction is \begin{equation}\label{source:eq-A10}
Z_N(\mathbf n)=
       \left(\bigotimes_{j=0}^{m-1}\mathcal L_N^{n_j}\right)S_N,
\end{equation} where \(\mathcal L_N^n\) is the original full-Hermite functional for
the signed power \(R^n\). Its crossing scalar and all first derivatives
at double roots are part of its definition. Conversion from \eqref{source:eq-A10} to the
specified zero-framed reduced invariant is the original unit-modulus
complete-color framing/component character. The separate full one-cut
modified convention retains \(i^{N+3}/N\); it also has zero exponential
rate.

These exact formulas follow from \cite[Sections 9--10]{FillingPaper},
applied to the displayed Morse diagram with the face order and parity
cochain above.

\section{The spectator envelope is exactly the union of the genuine
rates}\label{source:sec-3}

Use the source-normalized triangle potential \begin{equation}\label{source:eq-A11}
\begin{aligned}
Q(w_1,w_2,w_3)&=
 \max_{\substack{\alpha_i\ge0\\\sum_i\alpha_i=1/2}}
 \left\{\frac1\pi\sum_i
 [\Lambda(\pi\alpha_i)+\Lambda(\pi/2-\pi\alpha_i)]
       +\sum_i\alpha_i\log w_i\right\},\\
 \Lambda(t)&=-\int_0^t\log|2\sin u|\,du.
\end{aligned}
\end{equation} At a zero weight restrict its load to zero and omit its
logarithmic term from the sum. For one positive weight
it is one half of its logarithm. It is symmetric, nondecreasing, and
obeys \begin{equation}\label{source:eq-A12}
Q(cw)=Q(w)+\tfrac12\log c,
 \qquad Q(x,t,y)\ge\tfrac12\log t\quad(t>0).
\end{equation} The second inequality follows by using the admissible vertex with
load \(1/2\) on \(t\). Its entropy is zero since
\(\Lambda(0)=\Lambda(\pi/2)=0\). This proves the bound also when either
other coordinate is protected.

For a fixed nonnegative face vector with \(x+y>0\), the genuine uniform
local upper rates are \begin{equation}\label{source:eq-A13}
\begin{split}
 X&=Q_b+Q_c,\qquad Q_b=Q(x,b,y),\quad Q_c=Q(x,c,y),\\
 A_{\mathrm L}&=\max(\log c,Q_b+\tfrac12\log c)\quad(c>0),\\
 A_{\mathrm R}&=\max(\log b,Q_c+\tfrac12\log b)\quad(b>0).
\end{split}
\end{equation} An overloaded zero middle gives the zero one-sided polynomial and is
omitted. From \eqref{source:eq-A12}, \begin{equation}\label{source:eq-A14}
\begin{gathered}
 Q_b+\tfrac12\log c\le X,\qquad
 Q_c+\tfrac12\log b\le X,\\
 \max(X,A_{\mathrm L},A_{\mathrm R})
     =H:=\max(X,\{\log b:b>0\},\{\log c:c>0\}).
\end{gathered}
\end{equation} Thus the spectator \(H\) used by the three-vector proof is precisely
the maximum of the genuine local rate bounds. No mixed term has been
lost. This identity is already recorded in
\cite{SectorCover}, equation
(4), and in the final paragraph of the global LEFT theorem.


\begin{foundation}[Uniform local sector bounds]\label{found:envelopes}
Fix nonnegative face moduli \((x,b,c,y)\) with \(x+y>0\), keeping each
protected modulus exactly zero. Write \(H_{V,s}\) for the original
central or one-sided polynomial, where \(s\in\{\mathrm C,\mathrm L,\mathrm R\}\).
On the face torus with these moduli, its real-spectral upper exponential
rate is at most \(X,A_{\mathrm L},A_{\mathrm R}\), respectively, as
defined in \eqref{source:eq-A13}. The bound is uniform on the entire
spectral interval \([-2,2]\), includes both parities of \(N\), and
does not depend on the face phases. A one-sided term whose overloaded
middle has modulus zero is identically zero.
\end{foundation}

The proof is imported from \cite[Section 8]{Cofactors} and
\cite[Sections 1--4]{LeftEnvelope}. The following explanation records
which uniformity is required; the local recurrence and its global
transfer estimate are not being replaced by a positivity assumption.
Here ``uniform local upper rate'' means \begin{equation}\label{source:eq-A15}
\limsup_{N\to\infty}\frac1N
 \log\sup_{\substack{\sigma\in[-2,2]\\|L_f|=\rho_f}}
 |H_{V,s}(L;\sigma)|\le A_s(\rho),\qquad s=\mathrm C,\mathrm L,\mathrm R.
\end{equation} For CROSS this follows from the cleared-cofactor coefficientwise
real-interval interpolation bound and the positive complete triangle
factors. For LEFT the original omitted-half product and both unused
spectral products remain. The normalized interval continuant has
nonnegative face coefficients on \([-2,2\cos(\pi/N)]\); taking face
moduli therefore preserves the signed-continuant upper action. Its
complete endpoint optimization gives \eqref{source:eq-A13}. The remaining real cap has
width \(O(N^{-2})\). Since the actual spectral degree is at most \(d\),
polynomial exterior comparison costs only \begin{equation}\label{source:eq-A16}
\exp\left[d\,\operatorname{arcosh}
   \frac{6-2\cos(\pi/N)}{2+2\cos(\pi/N)}\right]=O(1).
\end{equation} This applies to a polynomial with complex coefficients at each fixed
face-phase tuple. RIGHT follows by the actual row/middle reversal, with
\(b,c\) exchanged in its rate. These are the proofs of
\cite{LeftEnvelope}, §§1--4,
and \cite{Cofactors}, §8. They include
both parity conventions and sampled spectral zeros. A polynomial
coefficientwise LEFT bound against its complete value is neither
asserted nor needed.

In this family the outer pole is always positive, so \(x+y>0\) at every
box and every one of the three proposed vectors. The zero-coordinate
limits in \eqref{source:eq-A11}--\eqref{source:eq-A15} are therefore within the proved domain.

\section{Complete packing and the three fixed shared
vectors}\label{source:sec-4}

Let \(\theta=\pi/m\). The complete packing has concentric pole
boundaries of radii \(a_0-1,a_0+1\), where \(a_0=\csc\theta\), and unit
rim disks centered at \(a_0e^{2ij\theta}\). The south disk is the
exterior of its boundary. Adjacent rim disks are tangent and all other
rim pairs are disjoint. Invert at the marked rim contact and multiply
lengths by two. Both poles then have curvature one, the two protected
faces become lines, and the remaining curvatures are \begin{equation}\label{source:eq-A17}
k_j=\frac{\cos\theta}{\sin^2\theta}
       [\cos\theta-\cos((2j-1)\theta)]
 =\frac{2\cos\theta}{\sin^2\theta}
              \sin(j\theta)\sin((j-1)\theta).
\end{equation} Thus \(k_0=k_1=0\), all other \(k_j>0\), and
\(k_2=k_{m-1}=4\cos^2\theta\).

The labeled simple-sphere packing and its positive cut polyhedron
identify the original augmented complement. The calibrated complete
CROSS action is \begin{equation}\label{source:eq-A18}
\Phi_m=2\sum_{j=0}^{m-1}Q(1,k_j,k_{j+1})
              -\sum_{j=2}^{m-1}\log k_j
        =\frac{\operatorname{Vol}(Y_m^\epsilon)}{2\pi}.
\end{equation} The equality uses the original marked triangle calibration. It is not
a full positive-endpoint source theorem. Deleting the two protected
faces leaves a connected ordinary graph, since both poles meet every
vertex of the surviving rim path. Thus its face Hessian has only the
common-scale kernel, as required by the Morse CROSS theorem. Every
matched edge has a positive outer and at least one nonmonomial triangle.

Put \begin{equation}\label{source:eq-A19}
a=\frac94,\quad C=\frac{27}4,\quad
 \lambda=\frac{a}{k_2}=\frac9{16\cos^2\theta},\quad
 w_j=\lambda k_j,
 \qquad L_{N_0}=L_{S_0}=1.
\end{equation} The whole vector \(W\) has \(w_2=w_{m-1}=a\) and \begin{equation}\label{source:eq-A20}
c:=w_3=w_{m-2}=a(4\cos^2\theta-1)\in[9/2,C].
\end{equation} Define \(W^{(1)}\) by replacing only \(w_2\) with 1, and \(W^{(0)}\)
by replacing only \(w_{m-1}\) with 1. Each shared occurrence of that
face is changed at once. Protected weights remain exactly zero.

The surviving one-sided anchors at \(j=0,1\) are LEFT; those at
\(j=2,m-1\) are RIGHT. Every other box has both possible one-sided
anchors. This follows directly from \eqref{braid:eq-2}, not from a reflected physical
coupon. Use \(W^{(1)}\) for \(1\mathrm L\), \(W^{(0)}\) for
\(0\mathrm L\), and \(W\) for all other anchors.

\section{Bounds for every even number of
syllables}\label{source:sec-5}

Write \(F(r)=Q(1,r,0)\). Scale only the ordinary rim coordinates by
\(t\in(0,1]\), retaining pole weights one, and let \(P(t)\) denote the
complete action at that vector. The pole angle in triangle \(j\) is \begin{equation}\label{source:eq-A21}
\alpha_j(t)=\arctan\frac1{\sqrt{t(k_j+k_{j+1})+t^2k_jk_{j+1}}},
 \qquad \alpha_0(t)=\frac\pi2.
\end{equation} The complete star gives \(\sum_j\alpha_j(1)=\pi\). For every other
triangle, with \(D^2=t(u+v)+t^2uv\), \begin{equation}\label{source:eq-A22}
-\frac{d\alpha}{d\log t}
 =\frac{t(u+v)+2t^2uv}{2D(1+D^2)}
 \le\sin\alpha\cos\alpha\le\alpha.
\end{equation} Hence \(\alpha_j(t)\le\alpha_j(1)/t\), and differentiation of the
entire coupled action gives \[
 \frac{dP}{d\log t}=2-\frac2\pi\sum_j\alpha_j(t)
       \ge1-\frac1t.
\] Integrating in the decreasing direction yields \begin{equation}\label{source:eq-A23}
P(\lambda)-\Phi_m\le\lambda^{-1}-1+\log\lambda
       \le D_0:=\frac{16}9-1+\log\frac9{16}<\frac{203}{1000}.
\end{equation} The bound is independent of the number of triangles; it uses their
original complete star, rather than summing independent local
distortions.

For \(3\le j\le m-2\), \eqref{source:eq-A17} gives \begin{equation}\label{source:eq-A24}
\frac13\le\frac{w_{j-1}}{w_j}\le2,
 \qquad \frac13\le\frac{w_{j+1}}{w_j}\le2.
\end{equation} The first ratio is \(\sin((j-2)\theta)/\sin(j\theta)\), increasing
with \(j\); the second is \(\sin((j+1)\theta)/\sin((j-1)\theta)\),
decreasing. Their extremal bounds use \(\sin3\theta\le3\sin\theta\) and
\(\sin4\theta/\sin2\theta=2\cos2\theta\le2\), valid also at \(m=6\).

The logarithmic derivative of \(Q\) in a positive weight lies in
\([0,1/2]\). Therefore \(F\) is increasing and \(F(r)-\log r\),
\(F(r)-\tfrac12\log r\) are decreasing. Drop the positive pole, scale by
the outer rim coordinate, and use \eqref{source:eq-A24}. At every interior box, \begin{equation}\label{source:eq-A25}
\begin{split}
 X_j-\log\max(w_{j-1},w_{j+1})
 &\ge F(1/3)+F(2)-\log2=:B_{\mathrm{int}}>43/100,\\
 Q(1,w_{j\pm1},w_j)-\tfrac12\log w_{j\pm1}
 &\ge F(2)-\tfrac12\log2=:C_{\mathrm{int}}>21/50.
\end{split}
\end{equation} At boxes 2 and \(m-1\), the middle weights are \(0,c\) and the outer
rim weight is \(a\). The same derivative bounds and \(c\le C\) give \begin{equation}\label{source:eq-A26}
\begin{split}
 X_j-\log c&\ge F(a)+Q(1,C,a)-\log C=:B_{\mathrm{near}}>7/25,\\
 Q(1,c,a)-\tfrac12\log c
       &\ge Q(1,C,a)-\tfrac12\log C=:C_{\mathrm{near}}>21/50.
\end{split}
\end{equation} At boxes 0 and 1, \(X=F(a)>\log a\). Thus every spectator has \(H=X\)
at \(W\). An anchored LEFT loss is
\(\min(X-\log c,Q_c-\tfrac12\log c)\), and the RIGHT formula is
symmetric. Equations \eqref{source:eq-A25}--\eqref{source:eq-A26}, less \eqref{source:eq-A23}, give a net loss greater
than \(7/100\) for every anchor other than \(0\mathrm L,1\mathrm L\). In
particular even the weakest coarse comparison gives
\(7/25-203/1000=77/1000>7/100\).

For \(1\mathrm L\), use \(W^{(1)}\). Its actual anchored rate is now
zero and its complete rate is \(F(1)\). The entire complete-action
change from \(W\) is \begin{equation}\label{source:eq-A27}
T(c)=2[F(1)-F(a)+Q(1,1,c)-Q(1,a,c)]+\log a.
\end{equation} It increases in \(c\): the logarithmic derivative at \(c\) of the
first \(Q\) is larger than that of the second. The only possible
spectator excess is at box 2, \begin{equation}\label{source:eq-A28}
E(c)=\max(0,\log c-F(1)-Q(1,1,c))\le E(C).
\end{equation} At box 3, put \(z=w_4\). Here \(z\ge1\), \(z/c\le2\), and
\(1/c\ge4/27\), so \begin{equation}\label{source:eq-A29}
X_3-\log\max(1,z)
 \ge F(4/27)+F(2)-\log2>8/25.
\end{equation} Every other box is unchanged. Therefore the anchored cost at
\(W^{(1)}\) is at most \begin{equation}\label{source:eq-A30}
P(\lambda)+T(C)+E(C)-F(1),
 \qquad S_0:=F(1)-T(C)-E(C)>37/100.
\end{equation} Its saving beyond \(\Phi_m\) exceeds \(37/100-203/1000=167/1000\).
Reflection of the scalar weight formulas by \(j\mapsto1-j\), accompanied
by exchanging pole labels, proves the same estimate for \(0\mathrm L\)
at \(W^{(0)}\). It does not change the ordered quantum data in \eqref{braid:eq-2}.

The fixed constants in these comparisons have the rational enclosures
proved in Appendix~\ref{app:constants}. The proof uses only finitely many
fixed analytic constants; the estimates above hold for every even \(m\ge6\).

We have proved, at every surviving anchor \((j,s)\), \begin{equation}\label{source:eq-A31}
A_{j,s}(L^{j,s})+\sum_{i\ne j}H_i(L^{j,s})
     -\sum_{f\,\mathrm{ordinary}}\log L_f^{j,s}
                        <\Phi_m-7/100.
\end{equation}

\section{From the three vectors to the full shared
remainder}\label{source:sec-6}

Expand \eqref{source:eq-A8} into genuine label strings in
\(\{\mathrm C,\mathrm L,\mathrm R\}^m\), omitting zero strings and the
all-CROSS string. Each remaining string contains a surviving one-sided
anchor. Assign it to one such anchor, and use that anchor's single whole
face vector from Section~\ref{source:sec-4}.

For fixed real spectra, apply the multivariate Cauchy coefficient
estimate to \eqref{source:eq-A7} on the ordinary face torus of that vector. The
protected variables are already zero and have no integration circle. The
local estimates \eqref{source:eq-A15}, and the exact spectator identity \eqref{source:eq-A14}, give \begin{equation}\label{source:eq-A32}
\frac1N\log\|D_{N,\text{string}}\|_{[-2,2]^m}
 \le A_{j,s}+\sum_{i\ne j}H_i
       -\frac dN\sum_{f\,\mathrm{ordinary}}\log L_f+o_m(1).
\end{equation} The same face is used everywhere with its same weight. There is no
independent local optimization, no exchange of a maximum and an infimum,
and no relaxation of the shared demands.

The target difference \(d/N-1=-1/N\) contributes only the fixed number
\(N^{-1}\sum_f\log L_f\). For fixed \(m\) there are finitely many
vectors and at most \(3^m-1\) strings. Thus summing their absolute
values costs \(\log(3^m-1)/N=o_m(1)\), not an extra exponential rate.
Equation \eqref{source:eq-A31} proves \begin{equation}\label{source:eq-A33}
\limsup_{N\to\infty}\frac1N\log\|D_N\|_{[-2,2]^m}
                    \le\Phi_m-7/100.
\end{equation} In particular, for every \(0<\eta<7/100\), the norm is at most
\(\exp[N(\Phi_m-\eta)+o_m(N)]\). This statement includes even and odd
levels. It requires no phase positivity of the shared source and no
strict gap for the all-CROSS term itself on this whole spectral cube.

\section{Full Hermite insertion introduces no new exponential
cost}\label{source:sec-7}

For \(H\) of degree at most \(d\), the exact finite filling theorem
states \begin{equation}\label{source:eq-A34}
|\mathcal L_N^n H|\le(1+|n|/2)\|H\|_{[-2,2]}.
\end{equation} Its proof is the original midpoint/full-Hermite formula with the
physical crossing multiplier, followed by sharp trigonometric Bernstein.
It retains value and derivative terms at double roots and the odd simple
endpoint. It applies to negative as well as positive powers, without a
square-root normalization or omitted Cartan factor.

Both \(D_N\) and \(S_{N,\mathrm C}\) have degree at most \(d\) in each
spectral variable. Applying \eqref{source:eq-A34} successively yields \begin{equation}\label{source:eq-A35}
|Z_N-Z_{N,\mathrm C}|
 \le\prod_{j=0}^{m-1}(1+|n_j|/2)\,
                     \|D_N\|_{[-2,2]^m}.
\end{equation} The polynomial need not factor and its shared targets are unchanged.
The stronger bound exploiting product-node vanishing is optional. For
each fixed tuple \(\mathbf n\), the product in \eqref{source:eq-A35} is independent of
\(N\); hence \eqref{source:eq-A33} is the actual filled scalar remainder rate. No
Fourier-tail argument is needed for this remainder.

\section{The nonzero original Morse contribution and its filled
volume}\label{source:sec-8}


\begin{foundation}[The nonzero geometric filling contribution]\label{found:period}
Consider the original general Morse diagram of a fixed connected simple
matched sphere, with its cyclic orders, crossing-disk frames and marked
strand edge. Assume its complete labeled packing is realized by the
original positive ideal-cell cut, the ordinary face graph is connected,
and each matched pair has a positive outer coordinate and at least one
nonmonomial triangle. Retain the original coefficient targets, integer
face lattice and full spectral numerator. Fix retained half twists
\(\epsilon_e\), and fill the original disks by
\(n_e=\epsilon_e-2p_e\). If the main strands form one knot, then for
all independently sufficiently large \(|p_e|\), the actual all-central
period has the nonzero asymptotic \eqref{source:eq-A36}. Its selected
action obeys \eqref{source:eq-A37}--\eqref{source:eq-A39}; in particular
it is the volume of the smoothly filled original knot exterior divided
by \(2\pi\), and tends to the complete augmented action as all filling
lengths tend to infinity. The main cusp remains complete. There is no
condition on the relative sizes or signs of the fixed filling integers.
\end{foundation}

The analytic conclusion is \cite[Sections 1--8]{MorsePeriod}; the
original topology, marked primitive words, smooth completion and absolute
volume calibration are \cite[Sections 1--9]{MorseGeometry}, reproduced
as the two parts of Appendix O in \cite{ProofCollection}. These are
substantial companion proofs. They give the central contribution, not
the full invariant. The missing complementary bound is proved in this
paper. The complete simple-sphere geometry used in their hypotheses is
constructed in \cite[Sections 2--3]{Packing}; the packing below also
provides its explicit circle configuration.

The simple sphere, positive labeled packing, connected ordinary graph,
primitive rank-\(m-1\) face quotient, and nonmonomial triangle at each
matching edge have been checked in Sections~\ref{braid:sec-1}--\ref{braid:sec-3} and~\ref{source:sec-4}. These are precisely the
hypotheses of \cite{MorsePeriod} and
\cite{MorseGeometry}. They concern the
original all-CROSS projection \eqref{source:eq-A7}; they do not require complete-radius
balance, a middle forest, a flat sweep or a full positive-endpoint
source theorem.

Choose retained \(\epsilon_j\in\{-1,0,1\}\) with the prescribed parities
and write \(n_j=\epsilon_j-2p_j\). Let \(A_N(\delta)\) be the finite
face-critical action in the companion period theorem: the shared face
variables have been evaluated at their selected stationary point,
with the original scalar removed as specified in that theorem. Define
\[
 \mathcal S_{N,\mathbf1}(\delta)
 =A_N(\delta)+\pi i\sum_j n_j\delta_j^2+2\pi i\sum_j\delta_j.
\]
Let \(\delta_{N,\mathbf1}\) be its selected critical point and
\(S_{N,\mathbf1}\) its critical value. Write \(H_N(\delta)\) for
the face Hessian on the primitive common-scale quotient, and
\(Q_{N,\mathbf1}=-\partial_\delta^2\mathcal S_{N,\mathbf1}
(\delta_{N,\mathbf1})\) for the negative spectral Hessian. The face
square root continues from the positive complete Hessian; the spectral
square root is its Gaussian branch. The raw period has the nonzero
asymptotic \begin{equation}\label{source:eq-A36}
\begin{split}
 Z_{N,\mathrm C}={}&q^{d^2\sum_jn_j/4}(-2i)^m\sqrt{2\pi}\,N^{m+1/2}\\
 &\times\frac{\prod_j\sin(2\pi\delta_{N,\mathbf1,j})}
 {\sqrt{\det H_N(\delta_{N,\mathbf1})}\sqrt{\det Q_{N,\mathbf1}}}
 e^{N S_{N,\mathbf1}}
 \left(1+O_{m,\mathbf n}(N^{-1})
           +O_{m,\mathbf n}(e^{-c_{m,\mathbf n}N})\right).
\end{split}
\end{equation} The companion theorem proves \(A_N\to A_0\) with every
fixed derivative, \(S_{N,\mathbf1}\to S_{\mathbf1}\), and convergence
of the coefficient remaining after the explicit action and level
factors are removed to a nonzero number. In particular that coefficient
is subexponential. The sine
factors and determinants do not vanish at the chosen fixed long filling.
The first reflected charges have already been paired by the exact finite
trace; nonzero amplitude does not follow merely from a formal even
phase. The source factor \((-1)^{dm}\) and Fourier factor
\((-1)^{(N+1)m}\) cancel exactly. Both level parities therefore have
this same nonzero normalization.

The limiting pure action \(A_0\) has \(A_0(0)=\Phi_m\), and the selected
first-charge action is \begin{equation}\label{source:eq-A37}
S_{\mathbf1}=A_0(\boldsymbol\delta)
           +\pi i\sum_jn_j\delta_j^2+2\pi i\sum_j\delta_j.
\end{equation} Let \(R_{0,j}\) be the similarity of the pure ordered
shaded-side map before inserting the retained half twist. The original
meridian threading follows its inverse, so its pure logarithm is
\(-\operatorname{Log}R_{0,j}\), as in
\cite[Section 5, equations (14)--(16)]{MorseGeometry}. Let
\(U_j,V_j\) be the continuous logarithmic similarities of
\(\mu_j,\lambda_j\), with their zero values at the complete augmented
structure fixing the branches. The actual primitive logs satisfy \begin{equation}\label{source:eq-A38}
U_j=-\operatorname{Log}R_{0,j}+2\pi i\epsilon_j\delta_j,
 \quad V_j=-4\pi i\delta_j,
 \quad U_j+p_jV_j=-2\pi i.
\end{equation} They are original disk words with the signed half-twist correction
once. The source differentiation and the marked regulator identity
cancel the peripheral boundary term in \eqref{source:eq-A37}. The nearby framed
representation, actual logarithmic link periods and the primitive
equation in \eqref{source:eq-A38} give a smooth filling of those very circles; the main
cusp remains complete. Consequently \begin{equation}\label{source:eq-A39}
2\pi\operatorname{Re}S_{\mathbf1}
       =\operatorname{Vol}(S^3\setminus K_{m,\mathbf n}),
 \qquad \operatorname{Re}S_{\mathbf1}\longrightarrow\Phi_m
 \quad\text{as all }|n_j|\to\infty.
\end{equation} At the selected charge \(\operatorname{Im}\delta_j>0\),
so \(\operatorname{Re}V_j>0\). The auxiliary longitude therefore has
nonunit modulus. The small logarithmic periods of its two primitive
generators are fixed before multiplying one by \(p_j\). Thus
\eqref{source:eq-A38} is the actual angular period \(-2\pi i\), not
only an equation after exponentiation. The companion metric completion
fills these original solid tori and keeps the main cusp complete. Its
absolute volume comparison uses the same ideal-cell chain and a
boundary homotopy whose volume tends to zero at the cusp cutoff.
This is the metric statement used in \eqref{source:eq-A39}; a regulator
identity alone would not supply it. The continuation permits
independently long signed powers, with no ratio condition.

Fix \(0<\eta<7/100\). Increase the long-filling threshold until both
Morse theorems apply and the action drop is less than \(\eta/2\).
Equations \eqref{source:eq-A33}, \eqref{source:eq-A35} and the nonzero asymptotic \eqref{source:eq-A36} give \begin{equation}\label{source:eq-A40}
\frac{Z_N-Z_{N,\mathrm C}}{Z_{N,\mathrm C}}
       =O_{m,\mathbf n}(e^{-\kappa_{m,\mathbf n}N})
       \quad\text{for some }\kappa_{m,\mathbf n}>0.
\end{equation} The unit-modulus framing character and the optional single modified
factor have no exponential rate. This proves \eqref{braid:eq-5}.

There are finitely many retained half-twist patterns for fixed \(m\).
Taking the maximum of their thresholds gives \(B_m\) for every
knot-producing parity, not only the illustrative pattern with two odd
boxes. Any signed odd \(n_j\) may, for example, be represented with
\(\epsilon_j=1\) and integer \(p_j=(1-n_j)/2\); positivity of \(n_j\) is
not presumed.


\appendix
\section{Exact enclosures for the fixed constants}\label{app:constants}

The inequalities in the proof depend on fixed arguments of \(Q\),
not on a sampled list of diagram sizes. For positive weights on their
support, put \(D^2=\sum_{i<j}w_iw_j\) and
\(a_i=\arctan(w_i/D)\). Strict concavity of the triangle entropy, or
direct differentiation with its sum constraint, gives
\begin{equation}\label{eq:Q-closed}
 Q(w)=\frac1\pi\sum_{w_i>0}
 \{\Lambda(a_i)+\Lambda(\pi/2-a_i)+a_i\log w_i\}.
\end{equation}
Indeed \(\sum a_i=\pi/2\), and these are the stationary loads
\(\alpha_i=a_i/\pi\) in \eqref{source:eq-A11}.
When the support consists of one weight, use \(Q(w)=\tfrac12\log w\)
directly. All arguments of \(Q\) in the following table are rational.
The symbols \(a,C,D_0,B_{\rm int},C_{\rm int},B_{\rm near},C_{\rm near},T,E,S_0\)
are those in \eqref{source:eq-A19}--\eqref{source:eq-A30}.

\begin{center}
\begin{tabular}{ll}
\toprule
Fixed quantity & Rational enclosure sufficient for the proof\\
\midrule
\(D_0\) & \(D_0<203/1000\)\\
\(F(a)-\log a\) & \(>0\)\\
\(B_{\rm int}\) & \(>43/100\)\\
\(C_{\rm int}\) & \(>21/50\)\\
\(B_{\rm near}\) & \(>7/25\)\\
\(C_{\rm near}\) & \(>21/50\)\\
\(F(4/27)+F(2)-\log2\) & \(>8/25\)\\
\(S_0=F(1)-T(C)-E(C)\) & \(>37/100\)\\
\bottomrule
\end{tabular}
\end{center}

Here is the enclosing arithmetic used to establish the table. Write
\(J(t)=\int_0^t\log(2\sin\pi u)\,du\). For \(0<t<1/2\), the Euler
sine product, followed by termwise integration of its positive
absolutely convergent series, gives
\begin{equation}\label{eq:J-series}
 J(t)=t\log(2\pi t)-t-
 \sum_{n=1}^{20}\frac{\zeta(2n)t^{2n+1}}{n(2n+1)}+R(t),
 \quad -\frac{2t^{43}}{21\cdot43(1-t^2)}\le R(t)\le0.
\end{equation}
The tail follows from \(\zeta(2n)<2\) and
\(n(2n+1)\ge21\cdot43\) for \(n\ge21\). The identity
\(\Lambda(\pi t)/\pi=-J(t)\) then evaluates
\eqref{eq:Q-closed} without quadrature.

Every interval endpoint is an integer multiple of \(10^{-40}\).
Addition is exact at this scale; products and quotients take the
minimum and maximum endpoint combinations and round outwards.
All divisors are first checked to exclude zero. Square roots use
integer floor and ceiling square roots. Machin's identity
\(\pi=16\arctan(1/5)-4\arctan(1/239)\) is enclosed by forty terms
of each alternating series and its next-term error. For \(n\ge2\),
\[
 \sum_{k=1}^{64}k^{-2n}\le\zeta(2n)
 \le\sum_{k=1}^{64}k^{-2n}
       +\frac{1}{(2n-1)64^{2n-1}},
 \qquad\zeta(2)=\pi^2/6.
\]
For logarithms, rescale by a power of two and use eighty terms of
\(2\sum_{j\ge0}z^{2j+1}/(2j+1)\), where
\(z=(x-1)/(x+1)\) and the enclosing interval satisfies \(0\le z<1/2\).
After eighty terms its remaining value is between zero and
\(2z^{161}/(161(1-z^2))\). For positive arctangents, iterate
\(\arctan x=2\arctan(x/(1+\sqrt{1+x^2}))\) until the upper endpoint
is at most \(1/2\), then use eighty alternating terms and the
next-term error. Every normalized angle in \eqref{eq:Q-closed} is
checked to lie strictly between zero and \(1/2\) before
\eqref{eq:J-series} is applied.

The portable verifier \cite{Constants} executes these rational rules
and asserts the displayed comparisons. It also checks that the argument
of the maximum defining \(E(C)\) is positive. Thus its selected branch
is certified rather than inferred from a decimal approximation.
The smallest certified value of
\(B_{\rm near}-D_0\) exceeds \(817/10000\), while the weaker
table already gives \(7/25-203/1000=77/1000>7/100\).
No optimizer or loop over \(m\) appears in this finite calculation.
The passage from these constants to every even \(m\ge6\) is the
trigonometric and action argument in Section~\ref{source:sec-5}.

\section*{Scope}
The theorem concerns the original closed braid and its original crossing
disks. No terminal tangles are inserted. The crossing powers may have
arbitrary signs, but their magnitudes must exceed a threshold depending
on the fixed number of syllables. The permutation test for one knot is
essential. The proof gives no growing-diagram or moving-color statement,
and makes no assertion that these knot types are disjoint from classes
that admit another Conway presentation. Its external proof inputs are
exactly the three companion results stated in the text.

\begin{thebibliography}{20}
\bibitem{Kashaev} R. M. Kashaev, \emph{The hyperbolic volume of knots
from the quantum dilogarithm}, Lett. Math. Phys. 39 (1997), 269--275.
\href{https://arxiv.org/abs/q-alg/9601025}{arXiv:q-alg/9601025}.
\bibitem{Murakami} H. Murakami and J. Murakami, \emph{The colored
Jones polynomials and the simplicial volume of a knot}, Acta Math.
186 (2001), 85--104. \href{https://arxiv.org/abs/math/9905075}{arXiv:math/9905075}.
\bibitem{FKP} D. Futer, E. Kalfagianni and J. S. Purcell,
\emph{Cusp areas of Farey manifolds and applications to knot theory},
Int. Math. Res. Not. 2010, 4434--4497.
\href{https://arxiv.org/abs/0808.2716}{arXiv:0808.2716}.
\bibitem{ProofCollection} \emph{Volume limits from planar tangles},
completed-proof collection, Appendices N and O.
\href{https://sashakolpakov.github.io/murakamis_dream/manuscript.pdf}{Proof collection}.
\bibitem{FillingPaper} \emph{Finite spectral evaluation and uniform filling
bounds}, Theorems 1.1, 9.1 and 10.1, and Corollary 9.2.
\href{https://sashakolpakov.github.io/murakamis_dream/papers/filling-bounds.pdf}{Companion paper}.
\bibitem{Cofactors} \emph{The actual crossing cofactor and its central
factorization}, Sections 1--3 and 8.
\href{https://sashakolpakov.github.io/murakamis_dream/notes/actual_cross_cofactor_factorization.html}{Cofactor proof}.
\bibitem{LeftEnvelope} \emph{The global LEFT weighted spectral envelope},
Sections 1--4.
\href{https://sashakolpakov.github.io/murakamis_dream/notes/actual_global_left_weighted_spectral_envelope.html}{One-sided envelope proof}.
\bibitem{SectorCover} \emph{Full-source volume from a cover of genuine
sectors}.
\href{https://sashakolpakov.github.io/murakamis_dream/notes/general_actual_sector_cover_filling_volume.html}{Shared-sector conventions}.
\bibitem{Packing} \emph{Balanced terminal decorations and volume limits
beyond tangle-operation trees}, Sections 2--3.
\href{https://sashakolpakov.github.io/murakamis_dream/notes/actual_balanced_terminal_fan_volume.html}{Original simple-sphere packing and cut}.
\bibitem{MorsePeriod} \emph{The geometric contribution and all Fourier
charges for matched spherical diagrams}, Sections 1--8.
\href{https://sashakolpakov.github.io/murakamis_dream/notes/actual_morse_cross_filling_period.html}{Central-period proof}.
\bibitem{MorseGeometry} \emph{The original complement and its geometric
contribution}, Sections 1--9.
\href{https://sashakolpakov.github.io/murakamis_dream/notes/actual_morse_geometric_contribution.html}{Primitive filling and absolute volume proof}.
\bibitem{Constants} \emph{Rational enclosures for the fixed constants
of the cyclic source proof}, portable verifier and rational output.
\href{https://github.com/sashakolpakov/murakamis_dream/blob/main/repro/check_cyclic_source_constants.py}{Verifier};
\href{https://github.com/sashakolpakov/murakamis_dream/blob/main/repro/results/cyclic-source-constants.json}{exact output}.
\end{thebibliography}
\end{document}
